% Options for packages loaded elsewhere
\PassOptionsToPackage{unicode}{hyperref}
\PassOptionsToPackage{hyphens}{url}
\documentclass[
]{article}
\usepackage{xcolor}
\usepackage[margin=1in]{geometry}
\usepackage{amsmath,amssymb}
\setcounter{secnumdepth}{-\maxdimen} % remove section numbering
\usepackage{iftex}
\ifPDFTeX
  \usepackage[T1]{fontenc}
  \usepackage[utf8]{inputenc}
  \usepackage{textcomp} % provide euro and other symbols
\else % if luatex or xetex
  \usepackage{unicode-math} % this also loads fontspec
  \defaultfontfeatures{Scale=MatchLowercase}
  \defaultfontfeatures[\rmfamily]{Ligatures=TeX,Scale=1}
\fi
\usepackage{lmodern}
\ifPDFTeX\else
  % xetex/luatex font selection
\fi
% Use upquote if available, for straight quotes in verbatim environments
\IfFileExists{upquote.sty}{\usepackage{upquote}}{}
\IfFileExists{microtype.sty}{% use microtype if available
  \usepackage[]{microtype}
  \UseMicrotypeSet[protrusion]{basicmath} % disable protrusion for tt fonts
}{}
\makeatletter
\@ifundefined{KOMAClassName}{% if non-KOMA class
  \IfFileExists{parskip.sty}{%
    \usepackage{parskip}
  }{% else
    \setlength{\parindent}{0pt}
    \setlength{\parskip}{6pt plus 2pt minus 1pt}}
}{% if KOMA class
  \KOMAoptions{parskip=half}}
\makeatother
\usepackage{color}
\usepackage{fancyvrb}
\newcommand{\VerbBar}{|}
\newcommand{\VERB}{\Verb[commandchars=\\\{\}]}
\DefineVerbatimEnvironment{Highlighting}{Verbatim}{commandchars=\\\{\}}
% Add ',fontsize=\small' for more characters per line
\usepackage{framed}
\definecolor{shadecolor}{RGB}{248,248,248}
\newenvironment{Shaded}{\begin{snugshade}}{\end{snugshade}}
\newcommand{\AlertTok}[1]{\textcolor[rgb]{0.94,0.16,0.16}{#1}}
\newcommand{\AnnotationTok}[1]{\textcolor[rgb]{0.56,0.35,0.01}{\textbf{\textit{#1}}}}
\newcommand{\AttributeTok}[1]{\textcolor[rgb]{0.13,0.29,0.53}{#1}}
\newcommand{\BaseNTok}[1]{\textcolor[rgb]{0.00,0.00,0.81}{#1}}
\newcommand{\BuiltInTok}[1]{#1}
\newcommand{\CharTok}[1]{\textcolor[rgb]{0.31,0.60,0.02}{#1}}
\newcommand{\CommentTok}[1]{\textcolor[rgb]{0.56,0.35,0.01}{\textit{#1}}}
\newcommand{\CommentVarTok}[1]{\textcolor[rgb]{0.56,0.35,0.01}{\textbf{\textit{#1}}}}
\newcommand{\ConstantTok}[1]{\textcolor[rgb]{0.56,0.35,0.01}{#1}}
\newcommand{\ControlFlowTok}[1]{\textcolor[rgb]{0.13,0.29,0.53}{\textbf{#1}}}
\newcommand{\DataTypeTok}[1]{\textcolor[rgb]{0.13,0.29,0.53}{#1}}
\newcommand{\DecValTok}[1]{\textcolor[rgb]{0.00,0.00,0.81}{#1}}
\newcommand{\DocumentationTok}[1]{\textcolor[rgb]{0.56,0.35,0.01}{\textbf{\textit{#1}}}}
\newcommand{\ErrorTok}[1]{\textcolor[rgb]{0.64,0.00,0.00}{\textbf{#1}}}
\newcommand{\ExtensionTok}[1]{#1}
\newcommand{\FloatTok}[1]{\textcolor[rgb]{0.00,0.00,0.81}{#1}}
\newcommand{\FunctionTok}[1]{\textcolor[rgb]{0.13,0.29,0.53}{\textbf{#1}}}
\newcommand{\ImportTok}[1]{#1}
\newcommand{\InformationTok}[1]{\textcolor[rgb]{0.56,0.35,0.01}{\textbf{\textit{#1}}}}
\newcommand{\KeywordTok}[1]{\textcolor[rgb]{0.13,0.29,0.53}{\textbf{#1}}}
\newcommand{\NormalTok}[1]{#1}
\newcommand{\OperatorTok}[1]{\textcolor[rgb]{0.81,0.36,0.00}{\textbf{#1}}}
\newcommand{\OtherTok}[1]{\textcolor[rgb]{0.56,0.35,0.01}{#1}}
\newcommand{\PreprocessorTok}[1]{\textcolor[rgb]{0.56,0.35,0.01}{\textit{#1}}}
\newcommand{\RegionMarkerTok}[1]{#1}
\newcommand{\SpecialCharTok}[1]{\textcolor[rgb]{0.81,0.36,0.00}{\textbf{#1}}}
\newcommand{\SpecialStringTok}[1]{\textcolor[rgb]{0.31,0.60,0.02}{#1}}
\newcommand{\StringTok}[1]{\textcolor[rgb]{0.31,0.60,0.02}{#1}}
\newcommand{\VariableTok}[1]{\textcolor[rgb]{0.00,0.00,0.00}{#1}}
\newcommand{\VerbatimStringTok}[1]{\textcolor[rgb]{0.31,0.60,0.02}{#1}}
\newcommand{\WarningTok}[1]{\textcolor[rgb]{0.56,0.35,0.01}{\textbf{\textit{#1}}}}
\usepackage{graphicx}
\makeatletter
\newsavebox\pandoc@box
\newcommand*\pandocbounded[1]{% scales image to fit in text height/width
  \sbox\pandoc@box{#1}%
  \Gscale@div\@tempa{\textheight}{\dimexpr\ht\pandoc@box+\dp\pandoc@box\relax}%
  \Gscale@div\@tempb{\linewidth}{\wd\pandoc@box}%
  \ifdim\@tempb\p@<\@tempa\p@\let\@tempa\@tempb\fi% select the smaller of both
  \ifdim\@tempa\p@<\p@\scalebox{\@tempa}{\usebox\pandoc@box}%
  \else\usebox{\pandoc@box}%
  \fi%
}
% Set default figure placement to htbp
\def\fps@figure{htbp}
\makeatother
\setlength{\emergencystretch}{3em} % prevent overfull lines
\providecommand{\tightlist}{%
  \setlength{\itemsep}{0pt}\setlength{\parskip}{0pt}}
\usepackage{bookmark}
\IfFileExists{xurl.sty}{\usepackage{xurl}}{} % add URL line breaks if available
\urlstyle{same}
\hypersetup{
  pdftitle={Ordinal and multinomial regression},
  hidelinks,
  pdfcreator={LaTeX via pandoc}}

\title{Ordinal and multinomial regression}
\author{}
\date{\vspace{-2.5em}}

\begin{document}
\maketitle

\textbf{Inference labs} use the five-step template: Hypothesis →
Visualise → Assumptions → Conduct → Conclude.

\subsection{Learning objectives}\label{learning-objectives}

\begin{itemize}
\tightlist
\item
  Fit a proportional-odds model with \texttt{MASS::polr()} and report a
  cumulative odds ratio.
\item
  Fit a multinomial logistic regression with \texttt{nnet::multinom()}
  for an unordered outcome.
\item
  Check the proportional-odds assumption and decide what to do when it
  fails.
\end{itemize}

\subsection{Prerequisites}\label{prerequisites}

Binary logistic regression.

\subsection{Background}\label{background}

Ordinal outcomes --- pain scales, tumour stages, Likert items --- carry
more information than a dichotomised version and usually less than a
continuous version. The proportional-odds model (\texttt{polr}) assumes
the effect of each predictor is the same across all cumulative splits of
the response; the single coefficient is the log cumulative odds ratio.
Multinomial logistic regression makes no such constraint and fits a
separate logit for each non-reference category.

The proportional-odds assumption is strong. A Brant test (in the
\texttt{brant} package) or simply comparing the \texttt{polr} fit to a
multinomial fit on deviance can detect violations. When it fails,
options include a partial proportional-odds model, the multinomial
model, or a continuation-ratio model.

Multinomial models are notoriously hungry for data: each non-reference
category requires a full set of coefficients. Reporting should list all
coefficient tables and a global test (likelihood ratio) rather than a
single \emph{p}-value.

\subsection{Setup}\label{setup}

\begin{Shaded}
\begin{Highlighting}[]
\FunctionTok{library}\NormalTok{(tidyverse)}
\FunctionTok{library}\NormalTok{(broom)}
\FunctionTok{library}\NormalTok{(MASS)}
\FunctionTok{library}\NormalTok{(nnet)}
\FunctionTok{set.seed}\NormalTok{(}\DecValTok{42}\NormalTok{)}
\FunctionTok{theme\_set}\NormalTok{(}\FunctionTok{theme\_minimal}\NormalTok{(}\AttributeTok{base\_size =} \DecValTok{12}\NormalTok{))}
\end{Highlighting}
\end{Shaded}

\subsection{1. Hypothesis}\label{hypothesis}

Simulate an ordinal outcome (symptom severity in three levels) driven by
a single continuous exposure, then contrast the ordinal and multinomial
analyses.

\subsection{2. Visualise}\label{visualise}

\begin{Shaded}
\begin{Highlighting}[]
\NormalTok{x }\OtherTok{\textless{}{-}} \FunctionTok{rnorm}\NormalTok{(n)}
\NormalTok{eta }\OtherTok{\textless{}{-}} \FloatTok{1.2} \SpecialCharTok{*}\NormalTok{ x}
\CommentTok{\# cumulative cutpoints}
\NormalTok{probs }\OtherTok{\textless{}{-}} \FunctionTok{cbind}\NormalTok{(}
  \FunctionTok{plogis}\NormalTok{(}\SpecialCharTok{{-}}\FloatTok{0.5} \SpecialCharTok{{-}}\NormalTok{ eta),}
  \FunctionTok{plogis}\NormalTok{(}\FloatTok{1.0} \SpecialCharTok{{-}}\NormalTok{ eta) }\SpecialCharTok{{-}} \FunctionTok{plogis}\NormalTok{(}\SpecialCharTok{{-}}\FloatTok{0.5} \SpecialCharTok{{-}}\NormalTok{ eta),}
  \DecValTok{1} \SpecialCharTok{{-}} \FunctionTok{plogis}\NormalTok{(}\FloatTok{1.0} \SpecialCharTok{{-}}\NormalTok{ eta)}
\NormalTok{)}
\NormalTok{y }\OtherTok{\textless{}{-}} \FunctionTok{apply}\NormalTok{(probs, }\DecValTok{1}\NormalTok{, }\ControlFlowTok{function}\NormalTok{(p) }\FunctionTok{sample}\NormalTok{(}\DecValTok{1}\SpecialCharTok{:}\DecValTok{3}\NormalTok{, }\DecValTok{1}\NormalTok{, }\AttributeTok{prob =}\NormalTok{ p))}
\NormalTok{dat }\OtherTok{\textless{}{-}} \FunctionTok{tibble}\NormalTok{(x, }\AttributeTok{y =} \FunctionTok{factor}\NormalTok{(y, }\AttributeTok{levels =} \DecValTok{1}\SpecialCharTok{:}\DecValTok{3}\NormalTok{,}
                            \AttributeTok{labels =} \FunctionTok{c}\NormalTok{(}\StringTok{"mild"}\NormalTok{, }\StringTok{"moderate"}\NormalTok{, }\StringTok{"severe"}\NormalTok{),}
                            \AttributeTok{ordered =} \ConstantTok{TRUE}\NormalTok{))}

\FunctionTok{ggplot}\NormalTok{(dat, }\FunctionTok{aes}\NormalTok{(y, x, }\AttributeTok{fill =}\NormalTok{ y)) }\SpecialCharTok{+}
  \FunctionTok{geom\_boxplot}\NormalTok{(}\AttributeTok{alpha =} \FloatTok{0.6}\NormalTok{, }\AttributeTok{colour =} \StringTok{"grey30"}\NormalTok{) }\SpecialCharTok{+}
  \FunctionTok{labs}\NormalTok{(}\AttributeTok{x =} \StringTok{"Severity"}\NormalTok{, }\AttributeTok{y =} \StringTok{"Exposure x"}\NormalTok{) }\SpecialCharTok{+}
  \FunctionTok{theme}\NormalTok{(}\AttributeTok{legend.position =} \StringTok{"none"}\NormalTok{)}
\end{Highlighting}
\end{Shaded}

\subsection{3. Assumptions}\label{assumptions}

Proportional-odds (same β across cumulative splits) for \texttt{polr};
independent observations for both models.

\begin{Shaded}
\begin{Highlighting}[]
\NormalTok{splits }\OtherTok{\textless{}{-}} \FunctionTok{tibble}\NormalTok{(}\AttributeTok{cut =} \FunctionTok{c}\NormalTok{(}\StringTok{"mild vs rest"}\NormalTok{, }\StringTok{"mild+mod vs severe"}\NormalTok{)) }\SpecialCharTok{|\textgreater{}}
  \FunctionTok{mutate}\NormalTok{(}\AttributeTok{coef =} \FunctionTok{c}\NormalTok{(}
    \FunctionTok{coef}\NormalTok{(}\FunctionTok{glm}\NormalTok{(}\FunctionTok{I}\NormalTok{(}\FunctionTok{as.numeric}\NormalTok{(y) }\SpecialCharTok{\textgreater{}=} \DecValTok{2}\NormalTok{) }\SpecialCharTok{\textasciitilde{}}\NormalTok{ x, }\AttributeTok{data =}\NormalTok{ dat, }\AttributeTok{family =}\NormalTok{ binomial))[}\DecValTok{2}\NormalTok{],}
    \FunctionTok{coef}\NormalTok{(}\FunctionTok{glm}\NormalTok{(}\FunctionTok{I}\NormalTok{(}\FunctionTok{as.numeric}\NormalTok{(y) }\SpecialCharTok{\textgreater{}=} \DecValTok{3}\NormalTok{) }\SpecialCharTok{\textasciitilde{}}\NormalTok{ x, }\AttributeTok{data =}\NormalTok{ dat, }\AttributeTok{family =}\NormalTok{ binomial))[}\DecValTok{2}\NormalTok{]}
\NormalTok{  ))}
\NormalTok{splits}
\end{Highlighting}
\end{Shaded}

If the two slopes are close, proportional odds is plausible.

\subsection{4. Conduct}\label{conduct}

\begin{Shaded}
\begin{Highlighting}[]
\FunctionTok{summary}\NormalTok{(fit\_polr)}
\CommentTok{\# approximate p{-}values}
\NormalTok{coef\_summary }\OtherTok{\textless{}{-}} \FunctionTok{coef}\NormalTok{(}\FunctionTok{summary}\NormalTok{(fit\_polr))}
\NormalTok{p }\OtherTok{\textless{}{-}} \FunctionTok{pnorm}\NormalTok{(}\FunctionTok{abs}\NormalTok{(coef\_summary[, }\StringTok{"t value"}\NormalTok{]), }\AttributeTok{lower.tail =} \ConstantTok{FALSE}\NormalTok{) }\SpecialCharTok{*} \DecValTok{2}
\FunctionTok{cbind}\NormalTok{(coef\_summary, }\StringTok{"p value"} \OtherTok{=} \FunctionTok{round}\NormalTok{(p, }\DecValTok{3}\NormalTok{))}
\FunctionTok{exp}\NormalTok{(}\FunctionTok{coef}\NormalTok{(fit\_polr))  }\CommentTok{\# cumulative odds ratio}
\end{Highlighting}
\end{Shaded}

Multinomial:

\begin{Shaded}
\begin{Highlighting}[]
\FunctionTok{summary}\NormalTok{(fit\_multi)}
\NormalTok{z }\OtherTok{\textless{}{-}} \FunctionTok{summary}\NormalTok{(fit\_multi)}\SpecialCharTok{$}\NormalTok{coefficients }\SpecialCharTok{/} \FunctionTok{summary}\NormalTok{(fit\_multi)}\SpecialCharTok{$}\NormalTok{standard.errors}
\NormalTok{p\_multi }\OtherTok{\textless{}{-}}\NormalTok{ (}\DecValTok{1} \SpecialCharTok{{-}} \FunctionTok{pnorm}\NormalTok{(}\FunctionTok{abs}\NormalTok{(z), }\DecValTok{0}\NormalTok{, }\DecValTok{1}\NormalTok{)) }\SpecialCharTok{*} \DecValTok{2}
\FunctionTok{round}\NormalTok{(p\_multi, }\DecValTok{3}\NormalTok{)}
\end{Highlighting}
\end{Shaded}

\subsection{5. Concluding statement}\label{concluding-statement}

\begin{quote}
The proportional-odds model returned a cumulative odds ratio per unit of
x of \texttt{round(exp(coef(fit\_polr)){[}1{]},\ 2)}, meaning larger x
shifted symptom severity upward. The multinomial model returned two
logits (moderate vs mild and severe vs mild) with effects in the same
direction; deviance-based comparison did not reject the
proportional-odds simplification.
\end{quote}

\subsection{Common pitfalls}\label{common-pitfalls}

\begin{itemize}
\tightlist
\item
  Reporting only the odds ratio and ignoring whether proportional-odds
  holds.
\item
  Using a multinomial model with a tiny reference category.
\item
  Treating an ordered outcome as continuous and fitting a linear model.
\end{itemize}

\subsection{Further reading}\label{further-reading}

\begin{itemize}
\tightlist
\item
  Agresti A. \emph{Analysis of Ordinal Categorical Data}.
\item
  Harrell FE. \emph{Regression Modeling Strategies}, ch.~13.
\item
  Venables WN, Ripley BD. \emph{Modern Applied Statistics with S},
  ch.~7.
\end{itemize}

\subsection{Session info}\label{session-info}

\begin{Shaded}
\begin{Highlighting}[]

\end{Highlighting}
\end{Shaded}

\subsection{See also --- chapter index}\label{see-also-chapter-index}

\begin{itemize}
\tightlist
\item
  \href{../index.Rmd}{Methods in this chapter}
\item
  \href{../../../ch00-find-your-method.Rmd}{Find your method (Chapter
  0)}
\end{itemize}

\end{document}
